
\chapter{Catene di Markov (Spazio Discreto)}

Le Catene di Markov sono una classe fondamentale di processi stocastici utilizzati per modellizzare fenomeni in cui l'evoluzione futura del sistema dipende unicamente dal suo stato presente, e non da come ci sia arrivato. Questa proprietà è nota come "assenza di memoria".

\section{Introduzione Euristica}

\begin{esempio}[Il meteo di Salisburgo]
A Salisburgo il meteo può essere Soleggiato (S), Piovoso (P) o Nevoso (N). Supponiamo valgano queste regole empiriche:
\begin{itemize}
    \item Non ci sono mai due giorni di sole di fila. Se oggi c'è sole, domani nevica o piove (50\% e 50\%).
    \item Se oggi piove o nevica, domani c'è una probabilità del 50\% che il tempo resti invariato, e 25\% per ciascuna delle altre due opzioni.
\end{itemize}
L'informazione su cosa farà domani è completamente determinata dal tempo di \emph{oggi}. Il fatto che ieri piovesse non aggiunge alcuna informazione utile. Questa è l'essenza della \textbf{proprietà di Markov}.
Le domande tipiche a cui risponde la teoria sono: "Qual è la probabilità che tra due settimane piova?" oppure "In media, quanti giorni di sole ci sono in un anno (comportamento a lungo termine)?"
\end{esempio}

\section{Definizione e Costruzione}

Sia $(\Omega, \F, \PP)$ uno spazio di probabilità e $(\F_n)_{n \in \N}$ la filtrazione di riferimento. Un processo stocastico $X = (X_n)_{n \in \N}$ a valori in uno spazio \emph{discreto} $E$ è una \textbf{Catena di Markov omogenea} se:
\[ \PP(X_{n+1} = y \mid \F_n) = \PP(X_{n+1} = y \mid X_n) = p(X_n, y) \quad \text{q.c.} \]
per ogni $y \in E$.

Gli ingredienti fondamentali per definire la dinamica del sistema sono due:
\begin{enumerate}
    \item La \textbf{distribuzione iniziale} $\mu_0$: la legge di $X_0$ (es. partiamo dallo stato $x$ con probabilità 1).
    \item Le \textbf{probabilità di transizione} $p(x, y)$: la probabilità di passare dallo stato $x$ allo stato $y$ in un passo. 
\end{enumerate}

Se lo spazio $E$ è finito ($|E| = N$), le probabilità di transizione possono essere racchiuse in una \textbf{Matrice di transizione} $M \in \R^{N \times N}$, in cui ogni riga somma a 1 (matrice stocastica):
\[ \sum_{y \in E} p(x, y) = 1 \quad \forall x \in E \]

\begin{theorem}[Esistenza e Unicità]
Dato uno spazio $E$, una distribuzione iniziale $\mu_0$ e un nucleo di transizione $p(x, y)$, esiste una Catena di Markov con tali caratteristiche, e la sua legge (cioè le probabilità su tutte le possibili traiettorie) è univocamente determinata.
\end{theorem}

\subsection{Dinamica a $n$ passi (Chapman-Kolmogorov)}

Se $p(x,y)$ descrive la mossa in 1 passo, indichiamo con $p^{(n)}(x,y)$ la probabilità di andare da $x$ a $y$ in $n$ passi. Sfruttando la proprietà di Markov, si ottiene l'equazione di \textbf{Chapman-Kolmogorov}:
\[ p^{(n+m)}(x, y) = \sum_{z \in E} p^{(m)}(x, z) p^{(n)}(z, y) \]
In forma matriciale (per $E$ finito), questo significa semplicemente che la matrice di transizione a $n$ passi è la potenza $n$-esima della matrice a 1 passo: $M^{(n)} = M^n$.

Un modo equivalente per guardare la dinamica è la \textbf{transizione di legge}. Se $\mu_n$ è il vettore colonna (o riga a seconda della convenzione) delle probabilità di trovarsi in ciascuno stato al tempo $n$, allora:
\[ \mu_{n+1} = \mu_n M \implies \mu_n = \mu_0 M^n \]

\section{La Proprietà di Markov Forte}

Un risultato molto profondo asserisce che la catena di Markov perde la memoria non solo agli istanti di tempo deterministici (es. $n=10$), ma anche agli istanti casuali descritti da \textbf{tempi d'arresto} $\tau$.

\begin{theorem}[Proprietà di Markov Forte]
Sia $\tau$ un tempo d'arresto. Condizionatamente all'evento $\{\tau < \infty\}$, il processo post-arresto:
\[ Y_k := X_{\tau + k} \]
è ancora una Catena di Markov con la \emph{stessa} matrice di transizione di $X$, indipendente dal passato osservato fino a $\tau$.
\end{theorem}

\section{Classificazione degli Stati}

\subsection{Comunicazione e Classi Irriducibili}
Diciamo che:
\begin{itemize}
    \item $y$ è \textbf{accessibile} da $x$ ($x \to y$) se c'è una probabilità positiva di arrivarci in un certo numero di passi ($p^{(n)}(x,y) > 0$).
    \item $x$ e $y$ \textbf{comunicano} ($x \leftrightarrow y$) se $x \to y$ e $y \to x$.
\end{itemize}
La comunicazione è una relazione di equivalenza. Partiziona $E$ in \textbf{classi irriducibili}. Se tutto $E$ forma un'unica classe, la catena si dice \textbf{Irriducibile} (si può viaggiare da qualsiasi stato a qualsiasi altro stato).

\subsection{Ricorrenza e Transitorietà}
Fissato lo stato di partenza $X_0 = x$, studiamo il tempo (eventuale) di primo ritorno in $x$:
\[ \tau_x := \inf\{n \ge 1 : X_n = x\} \]
Definiamo la probabilità di \emph{ritorno in $x$ in tempo finito}:
\[ F(x, x) = \PP_x(\tau_x < \infty) \]

\begin{definition}[Ricorrenza e Transitorietà]
Uno stato $x$ è:
\begin{itemize}
    \item \textbf{Ricorrente} se $F(x, x) = 1$ (è assolutamente certo che la particella tornerà in $x$).
    \item \textbf{Transiente} se $F(x, x) < 1$ (c'è una probabilità positiva che la particella se ne vada per sempre senza mai più tornare in $x$).
\end{itemize}
\end{definition}

Sia $V_x = \sum_{n=1}^\infty \1_{\{X_n = x\}}$ il numero totale di visite ad $x$. Si dimostra che:
\begin{itemize}
    \item Se $x$ è transiente, $V_x < \infty$ q.c. La particella \textbf{abbandona definitivamente} $x$. Criterio analitico: $\sum_n p^{(n)}(x,x) < \infty$.
    \item Se $x$ è ricorrente, $V_x = \infty$ q.c. La particella visiterà $x$ \textbf{infinite volte}. Criterio analitico: $\sum_n p^{(n)}(x,x) = \infty$.
\end{itemize}

\begin{hint}[Classificazione e Contagio]
In uno spazio finito non possono essere tutti stati transienti (sennò la particella "uscirebbe" dallo spazio, assurdo). 
Inoltre, la ricorrenza è "contagiosa": se $x$ è ricorrente e $x \to y$, allora anche $y$ deve essere ricorrente. Per questo, in una classe irriducibile, \textbf{gli stati sono o tutti ricorrenti o tutti transienti}. Se la catena è irriducibile ed $E$ è finito, essa è per forza Ricorrente.
\end{hint}

\section{Misure Invarianti e Teorema Ergodico}

\begin{definition}[Distribuzione Invariante]
Una distribuzione di probabilità $\mu^*$ su $E$ è \textbf{invariante} o stazionaria se per ogni $y$:
\[ \sum_{x \in E} \mu^*(x) p(x,y) = \mu^*(y) \iff \mu^* = \mu^* M \]
In termini di algebra, $\mu^*$ è un autovettore sinistro della matrice di transizione relativo all'autovalore 1. Se il sistema parte distribuito come $\mu^*$, la distribuzione rimarrà $\mu^*$ per tutti gli istanti successivi (equilibrio dinamico).
\end{definition}

Se uno stato è transiente, il suo "peso" all'equilibrio è zero: $\mu^*(x) = 0$.
Affinché esista una distribuzione invariante, la catena (se irriducibile) deve non solo tornare (ricorrente), ma \emph{tornare in fretta}. Questo introduce l'ultima classificazione per gli stati ricorrenti:

\begin{itemize}
    \item \textbf{Null-Ricorrente}: Ritorna al $100\%$ ($F=1$), ma il tempo medio di attesa per il ritorno è infinito ($\E_x[\tau_x] = \infty$). \emph{Non ammette misura di probabilità invariante}.
    \item \textbf{Positivo-Ricorrente}: Il tempo medio di ritorno è finito ($\E_x[\tau_x] < \infty$). \emph{Esiste una probabilità invariante, ed è proporzionale a:}
    \[ \mu^*(x) = \frac{1}{\E_x[\tau_x]} \]
\end{itemize}

\begin{spiegazione}[Misure stazionarie e tempi di ritorno]
C'è un profondo legame tra frequenza e probabilità limite. Se in media impiego 10 giorni a tornare a casa, significa che a lungo termine spenderò a casa 1 giorno su 10 (il 10\% del tempo). Ecco perché la probabilità stazionaria $\mu^*(x)$ è l'inverso del tempo medio di primo ritorno!
\end{spiegazione}

\begin{theorem}[Teorema Ergodico]
Sia $X$ una catena irriducibile e positivo-ricorrente (dunque con probabilità invariante $\mu^*$). Allora per ogni funzione $f$ ragionevole (integrabile rispetto a $\mu^*$), la media nel \textbf{tempo} converge alla media nello \textbf{spazio} quasi certamente:
\[ \lim_{n \to \infty} \frac{1}{n} \sum_{k=1}^n f(X_k) = \sum_{x \in E} f(x) \mu^*(x) \quad \text{q.c.} \]
\end{theorem}

Il Teorema Ergodico generalizza la Legge dei Grandi Numeri (LLN). La LLN richiede indipendenza; il teorema ergodico permette "memoria", a patto che la catena si mescoli bene (positivo-ricorrente) finendo per dimenticare le condizioni iniziali.

\section{Esercizi Tipici}

Un classico esercizio di esame (o progetto) richiede di, dato un piccolo grafo o una matrice $M$:
\begin{enumerate}
    \item Disegnare il grafo delle transizioni e identificare le \textbf{classi irriducibili}.
    \item Stabilire chi è \textbf{Transiente} (chi prima o poi finisce in una trappola e non torna) e chi è \textbf{Ricorrente}.
    \item Se lo spazio è finito, la classe ricorrente è automaticamente positivo-ricorrente. Si imposta quindi il sistema $\mu^* M = \mu^*$ unito a $\sum \mu^*(x) = 1$ per trovare l'unica \textbf{Distribuzione Invariante}.
    \item Trovare i tempi medi di ritorno calcolando $\E_x[\tau_x] = 1 / \mu^*(x)$.
    \item Usare il Teorema Ergodico per stimare integrali temporali a lungo termine.
\end{enumerate}
